\documentclass[12pt]{article}
\usepackage[utf8]{inputenc}
\usepackage[T1]{fontenc}
\usepackage[margin=1in]{geometry}
\usepackage{setspace}
\usepackage{xcolor}
\usepackage[colorlinks=true, urlcolor=blue!60!black, linkcolor=blue!60!black]{hyperref}
\setlength{\parindent}{0em}
\setlength{\parskip}{0.6em}
\renewcommand{\baselinestretch}{1.0}
\begin{document}
\begin{center}
{\LARGE \textit{United States v. Vaello Madero}, 596 U.S. 159 (2022)\par}
\vspace{0.5em}
{\large Case Brief --- Two-Page Edition\par}
\end{center}

\section{Memory Jogger}
Congress may exclude Puerto Rico residents from SSI; the territory's general exemption from federal income taxes supplies a rational basis (\href{https://www.law.cornell.edu/supremecourt/text/20-303}{596 U.S. 159}).

\section{Detailed Case Facts}
SSI pays means-tested benefits to low-income aged, blind, and disabled residents of the States and D.C.; Puerto Rico residents receive a capped territorial program instead. Jose Luis Vaello Madero began receiving SSI in New York, moved back to Puerto Rico in 2013, and kept receiving payments until the SSA learned of his move and sued to recover \$28,081.

\section{Procedural History}
He defended on Fifth Amendment equal-protection grounds and won in the District of Puerto Rico and the First Circuit; the Supreme Court reversed.

\section{Judicial Votes}
8--1. Kavanaugh wrote for Roberts, Thomas, Breyer, Alito, Kagan, Gorsuch, and Barrett; Thomas and Gorsuch concurred; Sotomayor dissented.

\section{Holding}
Under the Territory Clause and rational-basis review, Congress may treat Puerto Rico residents differently in benefits programs, given the territory's distinct tax status.

\section{Analysis of Opinions}
Kavanaugh relied on \textit{Califano v.\ Torres} and \textit{Harris v.\ Rosario}, permitting Congress to weigh the general balance of benefits and burdens without dollar-for-dollar accounting, and warned that benefits symmetry would invite tax symmetry. Thomas questioned \textit{Bolling}'s reverse incorporation. Gorsuch called for overruling the \textit{Insular Cases}. Sotomayor argued that SSI recipients owe no income tax anywhere, so the tax rationale is irrational, and that Puerto Ricans pay billions in federal taxes yet lack representation.

\section{Examples: Future Applications}
\textbf{Same-Side:} excluding Guam and Puerto Rico from a new IRS-administered credit is upheld. \textbf{Opposite-Side:} excluding residents of a low-revenue State cannot rely on the Territory Clause. \textbf{Fence-Sitter:} excluding Puerto Rico from a benefit funded by a payroll tax its workers pay.

\section{Critique}
Critics say the decision applies rational basis to aggregate jurisdictions rather than individual fit, leaves Footnote~1 undertheorized, and shows that overruling the \textit{Insular Cases} would not have changed the result.

\section{Key Quotations}
\begin{quote}
``In devising tax and benefits programs, it is reasonable for Congress to take account of the general balance of benefits to and burdens on the residents of Puerto Rico. In doing so, Congress need not conduct a dollar-to-dollar comparison of how its tax and benefits programs apply in the States as compared to the Territories, either at the individual or collective level.''
\end{quote}
\noindent --- Justice Kavanaugh, Opinion of the Court, 596 U.S. at 165--66.

\end{document}
